\documentclass[11pt]{article}
\usepackage[margin=1in]{geometry}
\usepackage[T1]{fontenc}
\usepackage[utf8]{inputenc}
\usepackage{amsmath,amssymb,mathtools}
\usepackage{booktabs}
\usepackage{hyperref}
\usepackage{xcolor}
\usepackage{listings}

\hypersetup{
  colorlinks=true,
  linkcolor=blue!60!black,
  urlcolor=blue!60!black
}

\definecolor{codebg}{RGB}{248,248,248}
\definecolor{codeborder}{RGB}{210,210,210}
\definecolor{codecomment}{RGB}{80,120,80}
\definecolor{codekeyword}{RGB}{30,60,150}

\lstdefinestyle{reportcode}{
  backgroundcolor=\color{codebg},
  basicstyle=\ttfamily\small,
  breaklines=true,
  columns=fullflexible,
  commentstyle=\color{codecomment},
  frame=single,
  framerule=0.5pt,
  keywordstyle=\color{codekeyword},
  rulecolor=\color{codeborder},
  keepspaces=true,
  showstringspaces=false,
  tabsize=2
}

\lstset{style=reportcode, language=Python}

\setlength{\parindent}{0pt}
\setlength{\parskip}{0.5em}

\newcommand{\R}{\mathbb{R}}
\newcommand{\ysp}{y_{\mathrm{sp}}}
\newcommand{\uvec}{u}
\newcommand{\xvec}{x}
\newcommand{\Aaug}{A_{\mathrm{aug}}}
\newcommand{\Baug}{B_{\mathrm{aug}}}
\newcommand{\Caug}{C_{\mathrm{aug}}}
\newcommand{\Qout}{Q_{\mathrm{out}}}
\newcommand{\Rin}{R_{\mathrm{in}}}
\newcommand{\Hp}{H_p}
\newcommand{\Hc}{H_c}
\newcommand{\argmin}{\mathop{\mathrm{arg\,min}}}


\title{Horizon Adaptation In RL-Assisted MPC}
\date{2026-04-01}

\begin{document}
\maketitle

\section*{Method Goal}

This method learns when the offset-free MPC controller should use different
prediction and control horizons. The key idea is that the actor does not change
continuous inputs directly. Instead, a DQN policy selects one discrete horizon
recipe from a finite menu and the MPC is rebuilt around that choice.

\section*{Notebook Sources Covered}

\begin{itemize}
\item \texttt{RL\_assisted\_MPC\_horizons.ipynb}
\item \texttt{RL\_assisted\_MPC\_horizons\_dist.ipynb}
\item Shared helpers from \texttt{utils/helpers.py}, \texttt{DQN/dqn\_agent.py},
      and \texttt{Simulation/mpc.py}
\end{itemize}

\section*{Shared MPC Context}

All horizon experiments reuse the same baseline ingredients:

\begin{itemize}
\item a nonlinear polymer CSTR plant for closed-loop rollout,
\item a linear augmented model $(\Aaug, \Baug, \Caug)$ for offset-free MPC,
\item a scaled observer state $\hat{x}_{t}$,
\item scaled setpoint deviations and scaled input deviations.
\end{itemize}

The MPC objective implemented in \texttt{Simulation/mpc.py} is
\[
J = \sum_{k=1}^{\Hp} \| y_{t+k} - \ysp \|_{\Qout}^2
  + \sum_{k=1}^{\Hc} \| \Delta u_{t+k-1} \|_{\Rin}^2.
\]
The horizon agent changes $\Hp$ and $\Hc$ while leaving the optimizer form
itself unchanged.

\section*{State, Action, Reward, And Update Logic}

The RL state is built from the scaled observer state, scaled setpoint deviation,
and scaled current input deviation:
\[
s_t =
\begin{bmatrix}
\tilde{x}_{t} \\
\tilde{y}_{\mathrm{sp},t} \\
\tilde{u}_{t}^{\mathrm{dev}}
\end{bmatrix}.
\]

The action is discrete:
\[
a_t \in \{0, 1, \dots, |\mathcal{H}| - 1\},
\quad
\mathcal{H} = \{(\Hp,\Hc) \mid \Hc \le \Hp\}.
\]
The selected index is converted to a horizon pair through
\texttt{action\_to\_horizons}. The reward is still computed from tracking and
input-move quality after the MPC solve. In the notebook family, the shaped
reward is multiplied by $0.01$ before it is stored in the DQN replay buffer.

The DQN update follows a Double-DQN target:
\[
y_t = r_t + \gamma
Q_{\theta^-}\left(s_{t+1},
\arg\max_a Q_{\theta}(s_{t+1}, a)\right).
\]

\section*{Core Algorithm Flow}

\begin{enumerate}
\item Build a finite recipe table of valid $(\Hp,\Hc)$ pairs.
\item Form the RL state from the scaled observer state and current operating
      point.
\item Let the DQN choose a horizon index, usually only every
      decision interval.
\item Rebuild the MPC object if the selected horizon pair changed.
\item Solve the MPC optimization with the active $\Hp,\Hc$, apply the first
      move to the plant, and update the observer.
\item Store $(s_t, a_t, r_t, s_{t+1})$ in the DQN replay buffer and run
      \texttt{train\_step()} after warm start.
\end{enumerate}

\section*{Key Code Paths}

\subsection*{Reusable helper for horizon recipes}

\begin{lstlisting}
def build_horizon_recipes(predict_grid, control_grid):
    recipes = []
    for Hp in predict_grid:
        for Hc in control_grid:
            if Hc <= Hp:
                recipes.append((Hp, Hc))
    if not recipes:
        raise ValueError("No valid Horizon recipes found")
    return recipes

def action_to_horizons(horizon_recipes, a_idx):
    return horizon_recipes[a_idx]
\end{lstlisting}

\subsection*{Notebook-local supervisor logic}

\begin{lstlisting}
if i > WARM_START:
    if (i % DECISION_INTERVAL == 0) or (last_a is None):
        if not test:
            a_idx = agent.take_action(current_rl_state.astype(np.float32),
                                      eval_mode=test)
        else:
            a_idx = agent.act_eval(current_rl_state.astype(np.float32))
        Hp, Hc = action_to_horizons(h_recipes, a_idx)
        if (Hp, Hc) != (current_Hp, current_Hc):
            rebuild_mpc(Hp, Hc)
            current_Hp, current_Hc = Hp, Hc
        last_a = a_idx
else:
    a_idx = a_idx_mpc
    Hp, Hc = action_to_horizons(h_recipes, a_idx)
\end{lstlisting}

\subsection*{DQN action selection and training}

\begin{lstlisting}
@torch.no_grad()
def take_action(self, state, eval_mode=False):
    self.steps += 1
    eps = 0.0 if eval_mode else self.eps_schedule.value(self.steps)
    if random.random() < eps:
        return random.randrange(self.action_dim)
    s = torch.from_numpy(state).float().unsqueeze(0).to(self.device)
    q1 = self.online.q1_forward(s)
    return int(q1.argmax(dim=1).item())

def train_step(self):
    ...
    a_star = self.online.q1_forward(ns).argmax(dim=1)
    q_t = self.target.combined_forward(ns, mode=self.target_combine)
    q_next = q_t.gather(1, a_star.view(-1, 1)).squeeze(1)
    y = r + self.gamma * q_next * (1 - done)
\end{lstlisting}

\section*{Variants Covered}

\begin{itemize}
\item The nominal notebook studies horizon adaptation without injected
      operating-condition drift.
\item The disturbance notebook adds gradual changes to process quantities such
      as $Q_i$, $Q_s$, and $hA$ through the shared schedule generator.
\item Both variants keep the same DQN-based horizon abstraction and the same
      observer-based state pipeline.
\end{itemize}

\section*{Limitations, Caveats, And Open Questions}

\begin{itemize}
\item The true execution logic is notebook-local. The core supervisor is not
      fully lifted into a reusable module.
\item The horizon decision interval is controlled by notebook globals rather
      than a shared config layer.
\item Reward scaling by $0.01$ is specific to this notebook family and should
      be preserved if results are compared historically.
\item \texttt{Simulation/rl\_sim.py} is not the source of truth for this
      method; the notebook-local supervisor is.
\end{itemize}

\end{document}
